% © 2026 Ing. David Jaroš — CC BY-NC-ND 4.0
%
% This work is licensed under a Creative Commons Attribution-NonCommercial-NoDerivatives
% 4.0 International License (CC BY-NC-ND 4.0).
%
% License History: Earlier drafts (up to v0.3) were released under CC BY 4.0.
% From v0.4 onward, all material is released under CC BY-NC-ND 4.0 to protect
% the integrity of the theoretical work during ongoing academic development.
%
% See LICENSE.md for full license text.

% bootstrap_step_m1_conformal.tex
% Modular Bootstrap Step M1: Conformal invariance of S[Θ] on T²
%
% Goal: Show S[Θ] restricted to T² is a 2D CFT with c = 3.
% Status: IN PROGRESS — see MODULAR_BOOTSTRAP_K1_PLAN.md §4 Step M1
% Parent plan: MODULAR_BOOTSTRAP_K1_PLAN.md

% UBT-AI-PROVENANCE-BEGIN
% schema: ubt-ai-provenance/v1
% tier: C_working
% ai_assistance: disclosed
% human_review: risk-based
% editorial_responsibility: Ing. David Jaroš
% policy: ../../AI_PROVENANCE.md
% notice: Working material; exhaustive human review is not claimed.
% UBT-AI-PROVENANCE-END

\ifdefined\INCLUDEMODE
\else
  \documentclass[12pt,a4paper]{article}
  \usepackage[a4paper, margin=2.5cm]{geometry}
  \usepackage{amsmath, amssymb, amsthm}
  \usepackage{mathtools}
  \usepackage{hyperref}
  \usepackage{xcolor}
  \usepackage{mdframed}
  \usepackage{booktabs}
  \usepackage{tcolorbox}

  \newtheorem{theorem}{Theorem}[section]
  \newtheorem{lemma}[theorem]{Lemma}
  \newtheorem{proposition}[theorem]{Proposition}
  \newtheorem{corollary}[theorem]{Corollary}
  \theoremstyle{definition}
  \newtheorem{definition}[theorem]{Definition}
  \theoremstyle{remark}
  \newtheorem{remark}[theorem]{Remark}

  \newmdenv[backgroundcolor=yellow!20, linecolor=orange, linewidth=1.5pt]{gapbox}
  \newmdenv[backgroundcolor=green!15, linecolor=green!60!black, linewidth=1.5pt]{resultbox}
  \newmdenv[backgroundcolor=red!8,   linecolor=red!60,   linewidth=1.5pt]{warnbox}
  \newmdenv[backgroundcolor=blue!8,  linecolor=blue!50,  linewidth=1.5pt]{insightbox}

  \title{\textbf{Modular Bootstrap Step M1}\\
         \large Conformal Invariance of $S[\Theta]$ on $T^2$}
  \author{Ing.\ David Jaro\v{s}}
  \date{2026-04-28}

  \begin{document}
  \maketitle
\fi

%──────────────────────────────────────────────────────────────────────────────
\section{Purpose and Status}
\label{sec:m1:purpose}
%──────────────────────────────────────────────────────────────────────────────

\begin{tcolorbox}[colback=blue!4!white, colframe=blue!50!black,
                  title=Step M1 — Conformal Invariance of $S[\Theta]$ on $T^2$]
\begin{tabular}{ll}
Parent plan      & \texttt{MODULAR\_BOOTSTRAP\_K1\_PLAN.md} \\
Goal             & Verify $S[\Theta]|_{T^2}$ is a 2D CFT with $c = 3$ \\
Milestone        & 2026-05-05 \\
Status           & IN PROGRESS \\
Acceptance       & $T(z)T(w)$ OPE coefficient reproduces $c = 3$ \\
\end{tabular}
\end{tcolorbox}

\begin{warnbox}
\textbf{No-fit rule}: No constant in this document may be chosen to match
$\alpha$ or reproduce $k = 1$.  Every numerical result must follow from
UBT axioms or be explicitly marked [MC].
\end{warnbox}

%──────────────────────────────────────────────────────────────────────────────
\section{Setup: UBT Action on the $\psi$-Torus}
\label{sec:m1:setup}
%──────────────────────────────────────────────────────────────────────────────

The UBT fundamental field is $\Theta(q, \tau) \in \mathbb{C} \otimes \mathbb{H}$
with $\tau = t + i\psi$ and the imaginary-time circle
$\psi \sim \psi + 2\pi R_\psi$.  Restricting to the imaginary-time torus
$T^2 = S^1_t \times S^1_\psi$ (compactifying also real time with period
$\beta = 1/(k_B T_{\mathrm{phys}})$ for the purposes of the CFT analysis):

The action on $T^2$ for the imaginary-part sector $\operatorname{Im}(\mathbb{H})$
is a free massless Gaussian:
\begin{equation}
  S[\Theta]\big|_{T^2}
  = \frac{1}{4\pi \alpha'} \int_{T^2} \bigl|\partial \Theta_{\mathrm{Im}}\bigr|^2
    \,d^2 z,
  \label{eq:m1:action}
\end{equation}
where $\Theta_{\mathrm{Im}} = (\Theta^1, \Theta^2, \Theta^3)$ are the three
$\operatorname{Im}(\mathbb{H})$ components, $z = \sigma + i\tau$ is the
worldsheet coordinate, and $\alpha' = R_\psi^2$ at the T-duality self-dual
point $R_\psi = R_t$.

This is the action of \textbf{three free massless compact bosons}
$X^i \equiv \Theta^i$, $i = 1, 2, 3$, each with target-space radius
$R = R_\psi$.

%──────────────────────────────────────────────────────────────────────────────
\section{Stress-Energy Tensor and OPE}
\label{sec:m1:TT}
%──────────────────────────────────────────────────────────────────────────────

For a single free compact boson $X$ with action
$S = \frac{1}{4\pi\alpha'}\int |\partial X|^2 d^2 z$,
the holomorphic stress-energy tensor is:
\begin{equation}
  T(z) = -\frac{1}{\alpha'} :\!(\partial X)^2\!:
  \label{eq:m1:T_single}
\end{equation}

The OPE is standard:
\begin{equation}
  T(z)\, T(w) = \frac{c/2}{(z-w)^4}
              + \frac{2\,T(w)}{(z-w)^2}
              + \frac{\partial T(w)}{z - w}
              + \text{regular},
  \label{eq:m1:TT_ope}
\end{equation}
with central charge $c = 1$ per free compact boson.

\begin{proposition}[Central charge of $S[\Theta]|_{T^2}$]
\label{prop:m1:c3}
The action $S[\Theta]|_{T^2}$ in~\eqref{eq:m1:action} is a 2D CFT with total
central charge
\begin{equation}
  c = \dim_{\mathbb{R}}(\operatorname{Im}\mathbb{H}) \times 1 = 3 \times 1 = 3.
  \label{eq:m1:c3}
\end{equation}
\end{proposition}

\begin{proof}
The three components $\Theta^1, \Theta^2, \Theta^3$ are decoupled (the action
is diagonal in the $\operatorname{Im}(\mathbb{H})$ basis, since the quaternion
metric on $\operatorname{Im}(\mathbb{H}) \cong \mathbb{R}^3$ is Euclidean).
Each component $\Theta^i$ contributes $c_i = 1$ independently.
The total central charge is $c = c_1 + c_2 + c_3 = 3$.
\end{proof}

\begin{resultbox}
$S[\Theta]|_{T^2}$ is a 2D CFT with $c = 3$.
The stress-energy tensor OPE coefficient in~\eqref{eq:m1:TT_ope} equals $c = 3$.
\textbf{Circularity check}: this result uses only $\dim_{\mathbb{R}}(\operatorname{Im}\mathbb{H}) = 3$
(a UBT axiom; no $\alpha$, no $m_e$). \quad [L0/L1 CLEAN]
\end{resultbox}

%──────────────────────────────────────────────────────────────────────────────
\section{Mode Expansion and Operator Content}
\label{sec:m1:modes}
%──────────────────────────────────────────────────────────────────────────────

On $T^2$ with modular parameter $\tau_{\mathrm{mod}}$, each compact boson $X^i$
has the mode expansion:
\begin{equation}
  X^i(z, \bar{z}) = x_0^i
    + \frac{\alpha'}{2}(p^i_L \ln z + p^i_R \ln \bar{z})
    + i\sqrt{\frac{\alpha'}{2}} \sum_{n \neq 0} \frac{1}{n}
      \bigl(\alpha^i_n z^{-n} + \tilde\alpha^i_n \bar{z}^{-n}\bigr),
  \label{eq:m1:mode_expansion}
\end{equation}
where the left/right momenta satisfy the quantisation condition from the
compactification $X^i \sim X^i + 2\pi R_\psi$:
\begin{equation}
  p^i_L = \frac{m^i}{R_\psi} + \frac{n^i R_\psi}{\alpha'},
  \qquad
  p^i_R = \frac{m^i}{R_\psi} - \frac{n^i R_\psi}{\alpha'},
  \label{eq:m1:momenta}
\end{equation}
with $m^i \in \mathbb{Z}$ (Kaluza-Klein number) and $n^i \in \mathbb{Z}$
(winding number).

\paragraph{Self-dual radius.}
At the T-duality self-dual point $R_\psi = \sqrt{\alpha'}$:
\begin{equation}
  p^i_L = m^i + n^i, \quad p^i_R = m^i - n^i.
  \label{eq:m1:selfdual_momenta}
\end{equation}
At this radius, the partition function per boson is exactly $\vartheta_3(\tau_{\mathrm{mod}})$:
\begin{equation}
  Z_{\mathrm{boson}}(\tau_{\mathrm{mod}}) = \frac{\vartheta_3(\tau_{\mathrm{mod}})}{|\eta(\tau_{\mathrm{mod}})|^2},
\end{equation}
so that the full partition function of $S[\Theta]|_{T^2}$ is:
\begin{equation}
  \hat{Z}(\tau_{\mathrm{mod}}) = \frac{\vartheta_3(\tau_{\mathrm{mod}})^3}{|\eta(\tau_{\mathrm{mod}})|^6}.
  \label{eq:m1:partition_function}
\end{equation}
The numerator $\vartheta_3^3$ is the partition function of the zero modes
(KK + winding); the denominator $|\eta|^6$ is the oscillator contribution.
This is consistent with the result already established in
\texttt{canonical/alpha/alpha\_best\_route.tex} Gap G8.

%──────────────────────────────────────────────────────────────────────────────
\section{Step M1 Result and Handoff to Step M2}
\label{sec:m1:result}
%──────────────────────────────────────────────────────────────────────────────

\begin{resultbox}
\textbf{Step M1 result}: $S[\Theta]|_{T^2}$ is a 2D CFT of three free compact
bosons at the self-dual radius, with $c = 3$ and partition function
$\hat{Z}(\tau_{\mathrm{mod}}) = \vartheta_3^3 / |\eta|^6$.

The operator spectrum consists of vertex operators
$V_{m^i,n^i}(z) = e^{i(p^i_L X^i_L + p^i_R X^i_R)}$
with dimensions $h = \frac{\alpha'}{4}(p_L^i)^2$,
$\bar{h} = \frac{\alpha'}{4}(p_R^i)^2$.

\textbf{Step M2} will use this operator content to compute the 4-point
function $\langle V V V V \rangle$ and set up the crossing symmetry constraint.
\end{resultbox}

\begin{gapbox}
\textbf{Open sub-question for Step M2}: Does the UBT gauge-consistency
constraint (Dirac quantisation on $S^1_\psi$ from
\texttt{canonical/appendices/appendix\_alpha\_geometry.tex \S1}) restrict which
$(m^i, n^i)$ sectors are present?  If only $m = 0$ or only $n \neq 0$ sectors
contribute, the operator spectrum is truncated and may distinguish $k = 1$
from $k \to \infty$.  This question is addressed in Step M4.
\end{gapbox}

\ifdefined\INCLUDEMODE
\else
  \end{document}
\fi
